% D3RA-DNA manuscript in the ACM journal format (acmart, acmsmall: single column, 6.75 x 10 in pages).
% Build: pdflatex main, bibtex main, pdflatex main (twice).  The IEEE version is kept in ieee_version/.
\documentclass[acmsmall,screen]{acmart}
% Venue metadata (journal, volume, DOI, rights) is assigned by ACM after acceptance; until then the ACM
% reference block and the running headers that would show placeholder volume numbers are switched off.
\setcopyright{none}
\settopmatter{printacmref=false}
\renewcommand\footnotetextcopyrightpermission[1]{}
\usepackage{graphicx}
\usepackage{textcomp}
\usepackage{multirow}
\usepackage{array}
\usepackage{enumitem}
\usepackage{tikz}
\usetikzlibrary{arrows.meta,positioning,fit,shapes.geometric,calc}
\usepackage{adjustbox}
\usepackage{etoolbox}
% the tables were laid out for a wider text block; scale any table wider than the ACM text block to fit
\BeforeBeginEnvironment{tabular}{\begin{adjustbox}{max width=\linewidth}}
\AfterEndEnvironment{tabular}{\end{adjustbox}}
\newcommand{\E}{\mathbb{E}}
\newcommand{\Prob}{\mathbb{P}}
\graphicspath{{figures/}}
\setlength{\emergencystretch}{3em}
\input{tables/macros}

\begin{document}

\title[Verified Stopping Dominates Adaptive Allocation in Multi-Container DNA Storage]{Verified Stopping Dominates Adaptive Allocation in Multi-Container DNA Storage: A Computational Study on a Digital Twin and Real Read Streams}

\author{Sai Mouli}
\email{saimouli12345678@gmail.com}
\affiliation{%
  \institution{M S Ramaiah Institute of Technology}
  \city{Bengaluru}
  \country{India}}

\author{J Geetha}
\email{geetha@msrit.edu}
\affiliation{%
  \department{Department of Computer Science and Engineering}
  \institution{M S Ramaiah Institute of Technology}
  \city{Bengaluru}
  \country{India}}

\author{Pruthviraj M. Savanur}
\email{rajsavanur2003@gmail.com}
\affiliation{%
  \institution{M S Ramaiah Institute of Technology}
  \city{Bengaluru}
  \country{India}}

\renewcommand{\shortauthors}{Mouli, Geetha and Savanur}

\begin{abstract}
Retrieving an erasure-coded DNA archive spread over several separately handled containers involves two
decisions: when to stop sequencing, and which container to read next.  We find that the first is worth an
order of magnitude more than the second.  All results are computational: a DT4DDS digital twin with six
independently synthesised, aged and amplified containers, and replayed public Illumina runs.  A streaming
decoder that stops only when every stripe, its CRC-32 and the file's SHA-256 verify consumed 3.49k reads
per retrieval, 27--32\% fewer than the 4.80--5.15k reads a non-stopping plan must reserve for 99\%
success, and returned no wrong file in \SafeN{} retrievals.  Testing for completion after every read
instead of every 50 would save under 1\% more.  Adaptive container choice is worth far less.  With a
single-efficiency container model, model-based rules read 7--49\% more than round-robin when one
container was near-dead (information-greedy, 441\%); the penalty largely disappears under a re-selected
likelihood or a dropout-aware model.  Fixed read proportions estimated offline were 4.2\% cheaper than
round-robin, and no adaptive rule reliably lowered the 0.99-quantile.  On the baseline no adaptive rule
beat round-robin by more than a few per cent; larger gains required dropout-dominated scenarios or a
prior calibrated on held-out reads ($-$16\% per retrieval, repaid after about 17 retrievals).
\end{abstract}

% ACM Computing Classification System; to be checked with the ACM CCS tool (https://dl.acm.org/ccs)
\ccsdesc[500]{Information systems~Information storage systems}
\ccsdesc[300]{Hardware~Emerging technologies}
\ccsdesc[300]{Mathematics of computing~Probability and statistics}

\keywords{DNA data storage, adaptive sequencing, verified stopping, erasure coding, coverage depth, digital twin, Bayesian calibration, multi-armed bandits, reproducibility}

\maketitle
% plain page numbers until ACM assigns the journal, volume and article number
\fancypagestyle{acmplain}{\fancyhf{}\fancyfoot[C]{\thepage}\renewcommand{\headrulewidth}{0pt}}
\pagestyle{acmplain}\thispagestyle{acmplain}

%=====================================================================
\section{Introduction}
%=====================================================================
DNA has been demonstrated as an archival medium at kilobyte to hundreds-of-megabyte
scale~\cite{church2012next,goldman2013towards,grass2015robust,erlich2017fountain,organick2018random}.
Retrieval requires sequencing, which samples molecules at random: some strands are read
many times, others never, and reads carry substitution, insertion and deletion
errors~\cite{heckel2019characterization,gimpel2023digital}.  Systems therefore protect each
strand with an inner code and groups of strands with an outer erasure code, and sequence to a
fixed depth expected to recover enough strands~\cite{grass2015robust,organick2018random}.  The
number of reads needed for decoding---the \emph{coverage depth}---is now studied in its own
right~\cite{barlev2025cover,gruica2024combinatorial}.

We consider archives split across several physically separate \emph{containers} (vials,
capsules, pools) that are synthesised, stored and handled separately.  This is a working
assumption rather than established practice; random access so far addresses files within one
pool~\cite{bornholt2016dna,organick2018random}, and recent theory treats erasure-coded data
partitioned over containers with uniform sampling~\cite{kong2026coverage} and container
failures~\cite{levy2026expected}.  Containers age, are amplified and fail differently, and
these differences are unknown when a retrieval starts.  Two operational questions follow:
\emph{when} to stop sequencing, and \emph{which} container to sequence next.

\textbf{Problem.}  Given an erasure-coded file distributed over $M$ containers of unknown,
heterogeneous quality, some of which may have failed, sequence in batches and decide after
each batch which container to read and whether to stop, so that the output file is
byte-identical to the stored one with high probability while the number of reads is small.

\textbf{Findings.}  (i) \emph{Verified stopping consumes 27--32\% fewer reads than the capacity a
non-stopping plan must reserve.}  On the digital twin, round-robin with verified stopping consumed 3.49k
reads per retrieval against the 4.80--5.15k that a non-stopping plan must reserve for 99\% success.  It
needs no container model and never returned a wrong file.  A smaller safety margin does not close the
gap: a non-stopping budget of 1.2 times the median fixed-depth requirement completes only 85\%
of retrievals.  Testing for completion after every read rather than every 50 reads saves a further
0.7\%, so verified stopping in batches captures nearly all of the available gain.  On the two
real-read pools, stopping reduced reads relative to fixed depth by 14\% and 25\%.
(ii) \emph{Choosing which container to read is worth a few per cent on the baseline.}  Fixed read
proportions estimated offline were 4.2\% cheaper than round-robin.  At the likelihood chosen by a
corrected selection, the model-based rules were no better than round-robin; under the pre-registered
likelihood they were 7--49\% worse when a container was near-dead.  An ablation traces this penalty to a
controller that is slow to recognise a bad container.  Larger allocation gains appeared only in
dropout-dominated scenarios and with a prior calibrated on held-out reads ($-$16\% per retrieval, repaid
after about 17 retrievals).
(iii) \emph{No adaptive rule reliably lowered the 0.99-quantile}, whereas fixed proportions did.
(iv) \emph{Two zero-success results reflect limits of the archive design rather than of the policies.}
Two failed containers leave exactly $K$ fragments per stripe by construction, and under the
photolithographic synthesis model no policy or static plan in our design recovered the file.

\textbf{Scope.}  All work is computational; nothing was synthesised or sequenced for this paper.  We use
the term \emph{exact retrieval} only for our digital-twin pipeline, where the output file is checked
byte for byte against the stored file.  On replayed real reads, decodability is an oracle event and we
speak of \emph{recovery}.

\textbf{Contributions.}
\begin{enumerate}
\item \emph{A measurement of the value of verified stopping.}  A verification-gated streaming pipeline
(Section~\ref{sec:method}) is compared with non-stopping plans that reserve the 0.99-quantile or a
margin above the median, and with the same schedule tested after every read, with reads consumed
reported separately from capacity reserved (Section~\ref{sec:res-actual}).  Stopping when
decodable~\cite{lopez2019dna} and checksum-verified decoding~\cite{erlich2017fountain} are both known;
our contribution is to quantify their value for heterogeneous multi-container retrieval, end to end to
verified bytes.
\item \emph{A negative result for adaptive allocation, with its cause isolated.}  Model-based rules,
including Thompson sampling, UCB and learn-then-commit~\cite{thompson1933likelihood,auer2002finite,garivier2016explore},
and two model-free rules, successive elimination and a yield threshold, are compared with round-robin and
with fixed proportions (Section~\ref{sec:res-alloc}).  A $2\times2\times2$ ablation separates the
container model from the likelihood and the prior (Section~\ref{sec:res-misspec}), and a fluid
approximation bounds what any allocation can gain (Section~\ref{sec:disc-theory}).
\item \emph{A reproducible end-to-end evaluation with explicit uncertainty.}  The full chain from
DT4DDS reads~\cite{gimpel2023digital} of independently synthesised containers to verified bytes is run
under pre-registered protocols.  We report safety bounds per retrieval and per synthesis replicate, tail
metrics with cluster-bootstrap intervals, sensitivity to the prior, the consensus cap, six further archive
geometries and four further failure modes, and a cost model in which reads are not the only charge.  Every
table and figure is regenerated from the released artifact (Section~\ref{sec:repro}).
\end{enumerate}

%=====================================================================
\section{Background and Related Work}\label{sec:related}
%=====================================================================
\textbf{DNA storage systems and channels.}  Storage systems combine constrained encodings with
inner and outer codes~\cite{antkowiak2020low,goldman2013towards,grass2015robust,erlich2017fountain,organick2018random,press2020hedges}.
Coverage is over-dispersed, a few percent of strands drop out, and errors arise mostly in
synthesis and sequencing~\cite{heckel2019characterization,chen2020quantifying,gimpel2023digital}.
DT4DDS~\cite{gimpel2023digital} is a digital twin of the synthesis, aging, PCR and sequencing
steps whose error and bias parameters were fitted to experiments; we use it as the simulator of
our primary evidence.  Information-theoretic models treat the channel as sampling with
replacement followed by noise~\cite{shomorony2021dna,weinberger2022dna,lenz2020coding}.

\textbf{Coverage depth.}  For a sampler that cannot be controlled, MDS codes and uniform
sampling minimise the expected number of reads~\cite{barlev2025cover}.  Follow-ups study random
access~\cite{gruica2024combinatorial}, real-data sampling distributions~\cite{cao2025optimizing},
single-container coverage depth of structured codes~\cite{bertuzzo2025coverage}, and the
reliability of MDS-coded retrieval~\cite{kashanna2026reliability}.  Cao
\emph{et al.}~\cite{cao2026modeling} drop the uniform-sampling idealisation that most of this line
assumes: they model the PCR and sequencing channel as nonuniform, instantiate it log-normally, and
analyse the MDS coverage-depth problem for both noiseless and noisy channels, deriving the expected
minimum depth for complete decoding in the noiseless case and, for the noisy case, a
coupon-collector envelope and a critical-quantile refinement, with sequencing sample-size estimates
validated by Monte-Carlo simulation.  Their object is one pool: how deep to sequence it so that enough
\emph{strands} are recovered.  That is the quantity our controller consumes rather than derives --- our
consensus-success curve and over-dispersed coverage model play the role their channel model plays --- and
it is complementary to what we study: several independently handled containers of unknown and
heterogeneous quality, some of which may have failed; \emph{which} container to read next after each
batch; when to stop, decided by verification rather than by a depth computed in advance; and, in the
digital-twin experiments, byte-for-byte verification of the returned file.  We take nonuniform sampling,
coverage-depth analysis and the use of checksums as established by this literature, not as contributions
of ours.

\emph{Distributed coverage depth.}  Two 2026 papers are closest to our setting.  Kong \emph{et al.}~\cite{kong2026coverage} define the distributed coverage-depth
problem: the $n$ columns of a rank-$k$ generator matrix are split over $M$ containers by an ordered
partition, and in every round \emph{one} column is drawn uniformly at random, with replacement and
\emph{noiselessly}, from \emph{each} container; $\tau(G,\mathcal P)$ is the first round at which the
drawn columns span $\mathbb F_q^k$.  They derive an exact formula for $\E[\tau(G,\mathcal P)]$ for
arbitrary linear codes and partitions by $q$-binomial M\"obius inversion, prove that an $[n,k]_q$ MDS
code is optimal for every fixed partition (in rounds and in total reads), give the lower bound
$M\,\E[\tau]\ge M\lceil k/M\rceil$ with $\min_{M,\mathcal P}M\,\E[\tau]=k$ attained by a staircase
$k$-partition, characterise one-round-perfect layouts, separate the regimes in which distribution
genuinely lowers total reads (constant-size containers) from those in which it only buys parallelism
(containers of size $k^\varepsilon$), and prove that the $q$-ary simplex code uniquely minimises
single-container coverage depth, resolving a conjecture of Bertuzzo \emph{et al.}~\cite{bertuzzo2025coverage}.
Their model is deliberately idealised: sampling is uniform within and across containers, containers are
homogeneous and never fail, reads are noiseless, every round costs $M$ reads whether or not they are
needed (their Remark II.1), and the allocation is fixed in advance.  Their Section~V lists
``nonuniform, asynchronous, adaptive, and noisy sampling'' as open problem~4 and an asynchronous
variant with a total read budget as an open problem.  Our work is an empirical study of that gap:
heterogeneous, noisy, possibly failed containers, reads allocated one batch at a time, and stopping
decided by verification rather than by an oracle span condition.  It is complementary to theirs; we
prove nothing about optimal codes or layouts, and we use their regime distinction (parallelism versus
genuine read saving) to interpret our results.  Levy
\emph{et al.}~\cite{levy2026expected} study a different objective: reconstructing the contents of one
\emph{failed} container from the surviving ones, with scalar MDS and MDS array codes.  Full-message
retrieval, which we study, is still required after such a repair, and the two stopping rules differ
(a prescribed missing block versus the whole message).  Neither work allocates reads adaptively, uses
noisy reads, or decodes to bytes.  These works analyse \emph{fixed} sampling laws; we study a
\emph{controllable} sampler empirically.

\textbf{Adaptive sequencing and sequential decisions.}  Nanopore adaptive sampling selects
molecules in real time~\cite{loose2016real,payne2021readfish,weilguny2023dynamic}, and
DNA-storage readouts have been sequenced until sufficient coverage~\cite{lopez2019dna}.
Choosing a container resembles a multi-armed bandit with a combinatorial ($K$-of-$N$ per stripe)
reward; we use Thompson sampling~\cite{thompson1933likelihood,russo2018tutorial}, upper
confidence bounds~\cite{auer2002finite}, explore-then-commit~\cite{garivier2016explore} and an
information-greedy rule from Bayesian experimental design~\cite{chaloner1995bayesian} as
baselines (see~\cite{lattimore2020bandit}).

\textbf{Placement and consensus.}  Erasure-coded storage limits the fragments of a stripe per
failure domain~\cite{weil2006crush,cidon2013copysets,ford2010availability,nath2006subtleties,leong2012distributed}.
Strands are reconstructed from noisy reads by trace
reconstruction~\cite{batu2004reconstructing,srinivasavaradhan2021trellis,sabary2024reconstruction}
after clustering~\cite{rashtchian2017clustering} or learned models~\cite{barlev2025deep};
systematic nanopore errors are documented~\cite{delahaye2021sequencing,wick2019performance}.
Table~\ref{tab:literature_comparison} positions this work.

\textbf{Relation to prior work.}  Prior DNA-storage work already stops sequencing when the decoder has
enough information~\cite{lopez2019dna} and verifies decoded data with
checksums~\cite{erlich2017fountain}, and the coverage-depth literature analyses the expected number of
reads until decodability, including across
containers~\cite{barlev2025cover,kong2026coverage,levy2026expected}.  Stopping, verification and
distributed DNA read allocation are therefore not contributions of this paper.  We did not find in that
literature a measurement, on an end-to-end pipeline that decodes to verified bytes, of how much a
verification-gated stop is worth relative to the capacity a static plan must reserve when containers
are heterogeneous and may have failed, or of whether choosing \emph{which} container to read next with a
learned container model helps or hurts under those conditions.  We provide that measurement, the
mechanism behind the negative allocation result, and an artifact from which every number can be
regenerated.

%=====================================================================
\section{System Model}\label{sec:model}
%=====================================================================
\begin{figure*}[t]
\centering
\resizebox{\textwidth}{!}{\input{figures/arch_dna.tex}}
\caption{Architecture of the study.  \textbf{(1)}~A file is bundled with its SHA-256, split into $S{=}25$
stripes of $K{=}32$ data fragments and extended by a Cauchy Reed--Solomon code to $N{=}48$; each fragment
becomes one 136-nt strand (address, payload, inner RS parity, CRC-16) and the strands are placed
round-robin over $M{=}6$ containers with different handling histories.  \textbf{(2)}~Each container is an
independent DT4DDS synthesis, aging, PCR and sequencing run that yields its own read stream.
\textbf{(3)}~Retrieval reads the streams in batches of 50 reads and makes two decisions after every batch:
\emph{which} container to read next (allocation) and \emph{whether to stop} (verified stopping, which
requires every stripe to decode and the CRC-32 and SHA-256 checks to pass).  The stored file is used only
after a retrieval ends, to count wrong outputs and false accepts.}
\label{fig:arch}
\end{figure*}

Table~\ref{tab:notation} lists the notation.
\input{tables/notation}

\subsection{Storage}
A file is bundled, hashed (SHA-256) and split into $S$ stripes of $K$ data fragments of $P$
bytes.  A systematic Cauchy Reed--Solomon code~\cite{reed1960polynomial} over GF($2^8$) extends
each stripe to $N$ fragments, any $K$ of which suffice.  Each fragment becomes a strand carrying
an address, an address check byte, the payload and an inner Reed--Solomon code with four parity
bytes plus a CRC-16, mapped to DNA by a whitened 2-bit code.  Code parameters, a CRC-32 per
stripe and the SHA-256 are kept in a manifest outside the DNA.  Fragments are assigned
round-robin to $M$ containers.

\subsection{Sequencing and recovery model used by the controller}
Reads from container $j$ form a stream.  The controller's model (used only to rank containers
and to forecast; never to decide that a fragment is recovered) treats the reads of a stored
fragment as negative-binomial with mean $\theta_j\,R_j\,\bar\pi_j$ after $R_j$ reads, where
$\bar\pi_j$ is the prior per-fragment read share (yield over stored fragments) and $\theta_j$ an
unknown efficiency multiplier ($\theta_j{=}0$: failed container).  A fragment becomes decodable
after $t$ reads with probability $\varphi(t)$ (consensus-success curve), and a molecule is absent
with probability $d$ (dropout).  Integrating over abundance gives the probability $Q(\mu)$ that a
stored fragment is decodable at mean coverage $\mu$:
\begin{equation}
Q(\mu)=(1-d)\,\E_g\Big[\sum_{n\ge0}\mathrm{NB}(n;k,g\mu)\,\varphi(n)\Big],
\label{eq:Q}
\end{equation}
with a shared log-normal strand effect $g$ (Gauss--Hermite quadrature).  Every container starts
from the same generic literature prior (yield 0.9, NB shape $k{=}10$, dropout 0.5\%, a generic
consensus curve, failure probability $10^{-3}$); no calibration data of the evaluated system is
used in the primary experiments.

\subsection{Objective and metrics}
A policy chooses, after each batch of $b$ reads, a container and whether to stop.  We report the
following quantities separately:
\begin{enumerate}
\item \textbf{Mean reads over all attempts} $\bar r_{\text{all}}$: the average reads consumed per
retrieval \emph{attempt}, counting attempts that end in failure (which consume the whole stream).
This is what a pay-per-read user pays on average.
\item \textbf{Mean reads conditional on success} $\bar r_{\text{ok}}$: the average over attempts that
returned the verified file.  It is the right metric only when the success probability is the same for
the policies compared, and it says nothing about the cost of failures.
\item \textbf{Success probability} $p_{\text{ok}}$: the fraction of attempts that returned the
verified file within the stream.
\item \textbf{Tail quantiles} q95 and q99 of the consumed reads, with cluster-bootstrap intervals.
\item \textbf{Reserved capacity} $B_{0.99}$: the budget a plan must commit \emph{in advance} to
succeed on 99\% of retrievals, i.e.\ the q99 of its read requirement.  Where every retrieval succeeds,
$B_{0.99}$ coincides numerically with (4); the two differ in who pays it.  A user who reserves a
flowcell pays $B_{0.99}$ on every retrieval, whether or not the retrieval needed it, whereas a
pay-per-read user pays (1).
\end{enumerate}
Verified stopping replaces the reserved capacity $B_{0.99}$ with the consumed mean $\bar r_{\text{all}}$;
its saving is therefore a comparison between a policy that stops and a plan that does not.  Comparing
two policies that both stop is a separate, allocation-only question (Section~\ref{sec:res-actual}).  No
metric here is a monetary cost; a cost model with per-run and per-container charges is given in
Section~\ref{sec:res-cost}.

%=====================================================================
\section{Online-Adaptive Retrieval}\label{sec:method}
%=====================================================================
\subsection{Posterior over container efficiency}
The support of $\theta_j$ is $\Theta=\{0\}\cup\{\theta^{(1)},\dots,\theta^{(25)}\}$, log-spaced in
$[0.01,5]$, with prior $p_0(0)=f_j$ and $p_0(\theta^{(\ell)})\propto(1-f_j)\exp(-\frac12(\ln\theta^{(\ell)}/s)^2)$,
with $s=0.5$ in every policy run.  After $R_j$ reads with $U_j$ verified fragments among the $n_j$ stored in container
$j$, we use a beta-binomial likelihood with intra-container correlation $\rho_U$:
\begin{equation}
p(\theta\mid R_j,U_j)\propto p_0(\theta)\,\mathrm{BB}\big(U_j;\,n_j,\,Q(\theta R_j\bar\pi_j),\,\rho_U\big),
\label{eq:posterior}
\end{equation}
where $\mathrm{BB}(u;n,p,\rho)$ has mean $np$ and variance $np(1-p)(1+(n-1)\rho)$; $\rho_U{=}0$
recovers the binomial.  Recoveries of a container's fragments are correlated (shared synthesis
bias, dropout), so the binomial is over-confident; $\rho_U$ is chosen on disjoint calibration
replicates (Section~\ref{sec:calib-protocol}).  We call the controller \emph{online-adaptive},
not self-calibrating: it adapts to observations, but its uncertainty is not well calibrated
(Section~\ref{sec:res-calib}).

\subsection{Allocation rules}
For an unrecovered fragment $i$ in container $j$, the hazard of becoming decodable within the
next $b$ reads is
$h_{ij}=[Q(\tilde\theta_j(R_j{+}b)\bar\pi_j)-Q(\tilde\theta_jR_j\bar\pi_j)]/[1-Q(\tilde\theta_jR_j\bar\pi_j)]$,
where $\tilde\theta_j$ is chosen by the rule:
\begin{itemize}
\item \emph{Round-robin + verified stop}: reads every container in turn (no model).
\item \emph{Coverage-greedy}: $\tilde\theta_j$ = posterior mean; reads $\arg\max_j\sum_i h_{ij}$
over fragments of unfinished stripes.
\item \emph{Pivotal}: weights each fragment by a normal-approximation sensitivity of its
stripe's completion probability (details in Appendix~\ref{supp:alg}).
\item \emph{Thompson sampling}~\cite{thompson1933likelihood,russo2018tutorial}: $\tilde\theta_j$
drawn from the posterior; coverage score.
\item \emph{UCB}~\cite{auer2002finite}: $\tilde\theta_j$ = posterior 0.9-quantile.
\item \emph{Learn-then-commit}~\cite{garivier2016explore}: five round-robin batches per
container, then the posterior is held fixed and the coverage score used.
\item \emph{Info-greedy} (Bayesian experimental design~\cite{chaloner1995bayesian}): reads the
container with the largest posterior variance of $\ln\theta_j$ until all fall below $10^{-3}$,
then coverage-greedy.
\item \emph{Tail-aware}: $\tilde\theta_j$ = posterior 0.1-quantile (pessimistic) and the score
counts only the 10\% of stripes with the largest remaining need (the bottleneck stripes that
set the retrieval's length); a heuristic aimed at the q99.
\end{itemize}
All rules probe each container with one batch first.  Table~\ref{tab:algorithms} gives the
operational specification of the loop and of every rule --- initialisation, batch size, tie-breaking,
when the posterior is updated, the stopping test and how an exhausted stream is handled.  The
\emph{static} family reads fixed
proportions $w_j\propto e_j^{\alpha}$, $\alpha\in\{0,0.5,1,2\}$, until decodable, where $e_j$ is
the per-read recovery rate of container $j$ measured in hindsight on the same scenario:
$\alpha{=}0$ is uniform fixed depth, $\alpha{=}1$ equalises expected coverage, and the best
$\alpha$ is a hindsight-tuned (``oracle'') static plan.

\subsection{Physical interpretation of container switching}
\label{sec:ops}
Our controller picks a container after every batch of $b{=}50$ reads.  This does not assume that a
sequencer can be redirected between physical vials that often, and the experiments do not demonstrate a
physical sequencing controller.  They model an \emph{allocation policy over $M$ independent read
streams}: each container has its own prepared library and its own stream of reads, and the policy
decides from which stream the next batch is consumed.  Three physical arrangements realise this
abstraction at different costs.  (i) \emph{Multiplexed run}: all $M$
libraries are barcoded and loaded on one flow cell; reads arrive interleaved and ``choosing a
container'' means choosing which barcode's reads to keep processing.  This matches our model closely
but does not reduce the reads the instrument actually produces, so it saves analysis and, on
pay-per-read services, money, but not instrument time.  (ii) \emph{Streaming (nanopore)}: reads
arrive in real time, and adaptive sampling can reject molecules mid-read~\cite{loose2016real,payne2021readfish};
with one flow cell per container, stopping is directly actionable, and with barcoded libraries
selective sequencing approximates container choice.  (iii) \emph{Sequential runs}: one run per
container, where a switch costs a library preparation and a run setup.  Our primary cost metric (reads)
describes case~(i) with per-read billing and is optimistic for cases~(ii) and~(iii); the per-batch and
per-container charges in Section~\ref{sec:res-cost} bound how much that optimism matters.  Verified
stopping, which is our main recommendation, needs no switching mechanism; the decoder only has to
signal that the file is complete.


\subsection{Verified stopping}
A retrieval stops with an output file only when (i) every stripe holds $K$ fragments that each
passed inner RS decoding, CRC-16 and the address check, (ii) every stripe matches its CRC-32, and
(iii) the file matches its SHA-256.  If (ii) fails, one more fragment is requested for that
stripe and the faulty fragment is localised by leave-one-out decoding.  If the streams are
exhausted first, the retrieval reports failure; it never returns an unverified file.  An
early-abandonment (futility) rule is disabled in all runs.

\subsection{Completion forecast}
The controller also forecasts the probability of finishing within a remaining budget, analytically and by
posterior simulation.  No policy uses the forecast; it was pre-specified, it is worse than a constant
forecast at the base rate, and it is reported in Appendix~\ref{supp:forecast}.

%=====================================================================
\section{Evaluation Methodology}\label{sec:eval}
%=====================================================================
\subsection{Digital-twin end-to-end pipeline (primary evidence)}
We store a fixed 18{,}432-byte random file with $P{=}24$-byte payloads, $(K,N)=(32,48)$, over
$M{=}6$ containers.  The stripe count follows from what is actually encoded: the library bundles the
file with a 162-byte header (name, length, hash) into an 18{,}594-byte blob, and
$\lceil 18{,}594/(32\times24)\rceil=25$ stripes of 32 data fragments cover it, the last one
zero-padded (24 stripes would hold only 18{,}432 bytes, which is the payload alone).  This gives
$25\times48=1{,}200$ strands of 136\,nt, 200 per container, 8 fragments of every stripe per container,
and, at 60 simulated reads per stored strand, a maximum stream of $60\times1{,}200=72{,}000$ reads.
All counts reported are those produced by the code rather than rounded design targets.  For each
container we run
DT4DDS~1.1.0~\cite{gimpel2023digital} with its published default parameters: array synthesis
(Twist model; log-normal oligo counts, $\sigma{=}0.3$), sampling of a given number of copies
per strand, aging, high-fidelity PCR and Illumina iSeq100 single-end 150\,nt sequencing,
producing 60 reads per stored strand.  The six containers are \emph{independent} syntheses with
their own seeds and histories.  The reads, after adapter trimming, go through inventory
clustering against the expected address prefixes, consensus (at most 24 reads per cluster),
inner decoding, outer decoding with CRC-32 repair and SHA-256 verification, in batches of 50
reads.  Every container decodes its own clusters; nothing uses the stored strands to decide
recovery.  The stored file and payloads are used only after the retrieval to count wrong
outputs and false accepts (a wrong payload that passed the inner checks).

\subsection{Scenarios}
\emph{Baseline}: fresh; aged one half-life; aged two half-lives with 10 copies per strand; 30
PCR cycles; 45 PCR cycles with 20 copies; near-dead (two half-lives, 40 cycles, $\sigma{=}0.5$,
3 copies).  \emph{One/two failed}: the fresh (and the 30-cycle) container returns only
non-informative reads.  \emph{Homogeneous control}: six fresh containers.
\emph{Photolithographic}: the baseline histories with DT4DDS's photolithographic synthesis
model.  \emph{Severity sweep}: the baseline with the near-dead container's copy number set to 200, 30, 10, 3
or 1.  \emph{Additional failure modes}: half of one container's designs physically lost before storage;
a quarter lost in each of two containers; one container whose reads are 90\% not derived from stored
strands (library-preparation failure or severe contamination); another at 50\%; and a common-cause
scenario in which every container is aged two further half-lives and amplified harder.
\emph{Additional geometries}: 3 and 12 containers, $(32,40)$, $(32,64)$ and $(16,24)$ codes, and a file
four times larger.  With three containers the six conditions cannot all be kept; we retain one healthy,
one weak and one near-dead container, so a third of that archive is near-dead instead of a sixth, which
makes it a harder archive and not only a smaller one.

\subsection{Protocol, statistics and multiple comparisons}\label{sec:eval-stats}
Design, policies, metrics, statistics, pre-specified comparisons and the calibration rule were written to
\texttt{dt\_protocol.yaml} before the primary runs and were not changed afterwards; we call this the
\emph{primary configuration}, and Appendix~\ref{app:dev} summarises the deviations.  Each scenario uses
30 independent syntheses of all six containers $\times$ 10 read orders (300 retrievals per policy); seeds
are derived with CRC-32 of the scenario, replicate and container; retrieval is bit-for-bit reproducible
given the committed read streams, but DT4DDS itself is not seed-deterministic, so the streams are part of
the artifact (Section~\ref{sec:repro}).  Confidence intervals are cluster-bootstrap percentile intervals
over syntheses (2{,}000 resamples~\cite{efron1993introduction}); comparisons are paired over the same
syntheses and read orders; wrong-output and false-accept rates carry exact Clopper--Pearson
bounds~\cite{clopper1934use}.  Every analysis script draws its resamples from one shared scheme that
depends only on the number of syntheses, so the same statistic on the same retrievals has the same
interval in every table and in the text.

The analyses fall into three groups.  \emph{Confirmatory}: the four comparisons pre-specified in
\texttt{dt\_protocol.yaml} (each policy against the static plans in the primary scenarios, the
tail-aware rule against coverage-greedy at q99, the homogeneous control and the severity sweep), the
$\rho_U$ selection rule and the forecast evaluation.  \emph{Secondary}: the analyses of
\texttt{dt\_protocol\_rev4.yaml}, which were specified before they were run (the dropout model, the
ablation, the $\rho_U$ re-selection, prior sensitivity, further failure modes and geometries, consensus
caps, the extended and calibration-estimated static families, the extended replicate set, the cost model
and the calibration cost).  \emph{Exploratory}: analyses chosen after results were seen (the binomial
likelihood, the complete comparison at $\rho_U{=}0.03$, the reservation baselines and the baselines on
regenerated streams).  All are reported whatever their direction.  Intervals are per-comparison 95\%
intervals and are not adjusted for multiplicity.  The confirmatory family alone contains several dozen
policy-scenario comparisons, so about one in twenty intervals is expected to miss its true value by
chance; our conclusions rest on effects that recur across scenarios, likelihood settings and the replay,
not on any single interval near zero, and secondary and exploratory results are read as supporting rather
than confirming evidence.

\subsection{Calibration protocol}\label{sec:calib-protocol}
$\rho_U$ is chosen from $\{0, 0.003, 0.01, 0.03, 0.1\}$ on ten \emph{disjoint} baseline
replicates (seeds 1000--1009) as the value whose 90\% credible interval covers the target
$\theta^\ast_j$ closest to 90\% of the time.  $\theta^\ast_j$ is the multiplier that best fits the
container's full recovery curve under the generic model, and the selection used prior sd $s{=}1$
whereas the policy runs use $s{=}0.5$ (a mismatch identified after the protocol was fixed; see
Appendix~\ref{supp:calib}); intervals are evaluated at the read
count where the prior predicts 40\% recovery and at twice that.  The selected value is then fixed
as the default for every policy run, evaluated on the held-out replicates 0--29, and used on
the real-read replay.  The completion forecast is evaluated with a budget equal to the median
coverage-greedy cost (Appendix~\ref{supp:forecast}).

\subsection{Real read streams (secondary evidence)}
We replay public Illumina iSeq runs of two pools from ENA PRJEB65931~\cite{gimpel2023digital,gimpel2023prjeb65931}:
a Twist pool after 0/2/7 days of aging or 15/60/90 PCR cycles, and a Genscript pool (aging 0\,d A/B,
2\,d, 7\,d; PCR 15/30 cycles).  Each run emulates one container holding 250 stripes of $(32,48)$,
i.e.\ $250\times48=12{,}000$ designed strands.  The Twist design file contains exactly 12{,}000
strands (108\,nt), all of which are used.  The Genscript design file contains 12{,}472 strands
(102\,nt); we draw a seeded random subset of 12{,}000 and \emph{exclude} the remaining 472 designs
entirely, discarding the reads that map to them: of the $6\times900{,}000$ evaluation reads, 86.3\%
map to a design and 4.23\% of those mapped reads are discarded as excluded designs, leaving
4{,}461{,}804 usable reads.  The excluded strands serve no other purpose, and the subset is the same
in every condition and policy, so all Genscript comparisons use identical read material.  We use 200
read permutations and a paired bootstrap.

\textbf{Real and emulated components of the replay.}  The replay is \emph{semi-empirical}.
\emph{Taken from the sequencing data}: the reads themselves, their number, which designed strand each
read came from, the per-container yield and abundance profile that follows from those counts, the
fraction of reads that map to no design, and the order in which reads arrive (we permute the real
evaluation reads rather than generating arrivals from a model).  \emph{Supplied by us}: the
partitioning of one pool's runs into ``containers''; the erasure-coded layout of 250 stripes over those
containers, which was never physically synthesised; and \emph{the decision that a fragment has been
recovered}.  This decision is an \emph{oracle} event: a fragment counts as recovered when the consensus
of its first $t$ reads in a random permutation of the run equals the designed strand, with $t$ drawn
from per-strand consensus curves estimated from real reads.  Our inner Reed--Solomon code, the CRC-16,
the address check, the outer code, the CRC-32 and the SHA-256 \emph{are not executed} on this path, and
no bytes are compared.  The replay therefore supports no statement about wrong outputs, false accepts or
byte-exact retrieval; those results come only from the digital twin, where the whole decode chain runs.
Nor is the replay a physical demonstration of the system, since the containers are aliquot runs of one
synthesis batch.  It supports the comparison of \emph{allocation and stopping policies} over read
streams with real arrival statistics, and we use it only as secondary evidence on the ranking of
policies.

\textbf{Separation of calibration and evaluation reads.}  Each run is read once and split by position: the first
300{,}000 reads are the calibration split and the next 900{,}000 the evaluation split, so no read is in
both.  Everything that generates a recovery event in the evaluation comes from the evaluation split
alone: the read order is a permutation of the evaluation reads, and the number of reads a strand needs
is drawn against per-strand consensus curves.  Two beliefs are used, and they differ exactly in which
split they may see.  The \emph{generic} prior (all primary replay results) sees no data from these pools
at all: yield 0.9, NB shape 10, dropout 0.5\%, a generic consensus curve.  The \emph{calibrated}
belief sees only the calibration split, i.e.\
per-container read counts and pooled consensus statistics estimated from calibration reads.  Evaluation
data is used deliberately in one place only: the hindsight-tuned static plan, which is labelled as an
oracle baseline wherever it appears and serves as an upper bound on what a static allocation could
achieve.  Our first analysis of these pools did not keep the splits separate (per-strand consensus
curves were estimated on the evaluation reads and static weights tuned on them); it is replaced, kept
only as an upper bound in Appendix~\ref{app:leak}, and no claim in this paper rests on it.

%=====================================================================
\section{Results}\label{sec:results}
%=====================================================================
\subsection{Safety of verified stopping}
\input{tables/safety}
Table~\ref{tab:safety} reports the safety record together with the dependence structure behind it.  It
pools every digital-twin retrieval in this paper except two sets of re-runs on regenerated read streams
(the clean reproduction of Section~\ref{sec:repro} and the baselines of Section~\ref{sec:res-extra}),
which repeat replicate labels rather than adding independent ones.  Across \SafeN{} end-to-end retrievals
no output file differed from the stored file in a single byte, and in \SafeAtt{} fragment decode attempts
no wrong payload passed inner RS decoding, the CRC-16 and the address check.  Unsuccessful retrievals
(one or two failed containers) ended by exhausting the streams and reported failure instead of returning
a file.

These retrievals are \emph{not} independent experiments.  They come from \SafeUnits{} distinct synthesis
replicates, \SafePerUnit{} retrievals per replicate on average, and the decode attempts inside one
retrieval share the same reads.  A Clopper--Pearson bound that treats every retrieval as an independent
Bernoulli trial would therefore overstate the evidence.  Per retrieval the upper limit is \SafeWUB{};
treating one synthesis replicate as one unit, which is conservative because it discards the information
in repeated read orders, it is \SafeWUBc{}, two orders of magnitude weaker.  The false-accept limits are
\SafeFUB{} per attempt and \SafeFUBc{} per synthesis replicate.  Zero observed errors does not imply zero
error \emph{probability}, and the bounds apply only to \emph{this} archive configuration: this code and
rate, this placement, this decoder with its 24-read consensus cap, DT4DDS's error model, and the
scenarios we ran.  Mechanistically, a false accept would require a wrong payload to pass a 4-symbol RS
code, a CRC-16 and an address byte, then the stripe CRC-32 and the file SHA-256; the cryptographic hash is
the final safeguard, and no experiment of this size can bound \emph{its} failure rate.

\subsection{Reads consumed versus capacity reserved}\label{sec:res-actual}
\input{tables/r4_actual}
The largest effect in this paper is the difference between reads that are \emph{consumed} and capacity
that must be \emph{reserved} (Table~\ref{tab:actual}).  In the baseline digital twin, round-robin with
verified stopping consumed 3.49k reads on average (95\% CI 3.37--3.60k), while the same uniform read
proportions without stopping must reserve 5.15k to succeed on 99\% of retrievals, and the best static
plan in the family must reserve 4.80k.  Verified stopping therefore consumes 27--32\% fewer reads than
the capacity a non-stopping plan must reserve, depending on which plan the reservation is taken from, and
it needs no model of the containers.

\input{tables/r9_baselines}
Because the q99 is a demanding reference, Table~\ref{tab:r9base}(a) adds non-stopping plans with smaller
margins.  A fixed-depth budget of 1.2 times the median requirement (4.08k) completes only
85\% of retrievals (95\% CI 80--90\%), and 1.2 times the mean
(4.19k) completes 89\%.  If a failed attempt is repeated with a fresh budget,
these plans spend 4.78k and 4.69k reads per completed retrieval, against 3.49k with
verified stopping; the same holds for the coverage-equalising proportions.  A margin-based reservation is
therefore cheaper than the q99 but fails often, and neither matches stopping.

How much of the possible gain does stopping in batches of 50 reads capture?  Applying the stopping test
after every read, with the same container schedule, gives the first read count at which this decoder can
return the verified file.  On regenerated read streams (Table~\ref{tab:r9base}(b)), this per-read stop
used 0.7\% fewer reads than batched round-robin (95\% CI $-$0.7 to $-$0.6\%), so batching costs less
than one per cent.

Allocation is a separate and much smaller effect.  Three static plans are reported.  All three read fixed
proportions $w_j\propto e_j^{\alpha}$ \emph{with} the same verified stopping and are paired against
round-robin on the same 300 syntheses and read orders.
\begin{enumerate}
\item \emph{Pre-specified plan} (primary).  $\alpha\in\{0,0.5,1,2\}$ was the pre-specified grid and
$\alpha{=}1$, which equalises expected coverage, is its best member: 3.343k reads, $-$4.21\% against
round-robin's 3.490k (95\% CI $-$4.9 to $-$3.5\%).
\item \emph{Extended grid} (secondary).  Adding $\alpha\in\{0.25,0.75,1.25,1.5,3\}$ moves the point
estimate to $-$4.57\% at $\alpha{=}1.5$ (CI $-$5.5 to $-$3.6\%).  The intervals overlap almost entirely, so
the finer grid does \emph{not} establish a better plan; it shows that the family is flat near its
optimum ($-$3.7\% to $-$4.6\% for $\alpha\in[0.75,2]$, decaying to $-$0.70\% at $\alpha{=}3$).
\item \emph{Calibration-estimated plan} (secondary).  The same $\alpha{=}1$ form with the per-container
efficiencies $e_j$ measured on the ten \emph{disjoint} calibration replicates, so the plan sees no
evaluation data: 3.338k reads, $-$4.35\% (CI $-$5.1 to $-$3.6\%), statistically indistinguishable from the
pre-specified hindsight plan.  The static advantage is therefore not an artefact of hindsight.
\end{enumerate}

\subsection{What allocation is worth}\label{sec:res-alloc}\label{sec:res-lik}
\input{tables/r6_likelihood}
\input{tables/r4_ablation}
The primary allocation comparison runs at $\rho_U{=}0.1$, a value chosen by a procedure that used the
wrong prior width (Appendix~\ref{supp:calib}); the corrected procedure picks
$\rho_U{=}0.03$.  We therefore re-ran the complete comparison at $\rho_U{=}0.03$, with all policies, all
three primary scenarios and the full 30 syntheses $\times$ 10 read orders (Table~\ref{tab:r6lik}).
Round-robin never evaluates the posterior, so its retrievals are bit-identical across likelihood
settings and serve as a common reference.  At $\rho_U{=}0.03$, four of the seven model-based rules are
statistically indistinguishable from round-robin on the baseline: coverage-greedy $-$0.4\% ($-$2.3 to $+$1.6\%),
Thompson sampling $+$0.3\% ($-$1.7 to $+$2.5\%), UCB $+$1.5\% ($-$0.5 to $+$3.7\%) and tail-aware $-$1.2\% ($-$3.2 to $+$0.7\%).  Pivotal ($+$12.2\%),
learn-then-commit ($+$34.9\%) and information-greedy (Table~\ref{tab:suppinfo})
remain clearly worse.  The $2\times2\times2$ ablation (Table~\ref{tab:r4abl}) shows the same pattern from
another angle: the large penalty appears only in the primary corner (single-efficiency model,
$\rho_U{=}0.1$, prior width $s{=}0.5$).  Relaxing any one of the three removes most of it, leaving the
rules between 4\% below and 13\% above round-robin, and combinations reach a floor of about $-$4.5\%.  On
the baseline, therefore, no adaptive rule beat round-robin by more than a few per cent in any likelihood,
model or prior setting we tested, which is about what fixed proportions already give ($-$4.21\%).  In the
homogeneous control, where there is no near-dead container to misread, every rule in
Table~\ref{tab:r6lik} is between $-$3.7\% and $-$2.4\% against round-robin under both likelihood settings.

\input{tables/dt_main}
\begin{figure*}[t]
\centering
\includegraphics[width=\textwidth]{fig_dt_policies.pdf}
\caption{Digital-twin end-to-end retrieval, baseline unless stated.  \textbf{(a)} For each policy, the
reads a retrieval actually \emph{consumes}: filled circles are the mean over 300 retrievals with a 95\%
cluster-bootstrap interval over syntheses, open squares the 0.99-quantile of the same retrievals.  The two
vertical lines are the capacity a plan that does \emph{not} stop must \emph{reserve} to succeed on 99\% of
retrievals: dotted for uniform fixed depth, dashed for the best static plan.  The horizontal distance from a
filled circle to a vertical line is the saving verified stopping provides; the distance between circles is
the (much smaller) allocation effect.  \textbf{(b)} Allocation effect alone: each bar is the paired
difference in mean reads against round-robin with the same verified stopping, computed on the retrievals in
which both policies succeeded, with 95\% cluster-bootstrap intervals; negative is cheaper than round-robin.
Information-greedy is annotated rather than plotted because it runs off the scale.  Scenarios in which no
retrieval succeeded (photolithographic, two failed containers) have no paired statistic and are omitted.}
\label{fig:dt}
\end{figure*}
Under the pre-registered configuration the penalty is large (Table~\ref{tab:dt}, Fig.~\ref{fig:dt}).  In
the baseline scenario every policy recovered the file in all 300 retrievals.  Every model-based rule was
costlier than round-robin on the same syntheses and read orders: Thompson sampling by 7\% (4--11\%),
tail-aware by 13\% (9--18\%), coverage-greedy by 15\% (10--21\%), UCB by 15\% (11--21\%), learn-then-commit by
45\% (40--49\%) and pivotal by 49\% (39--60\%).  The ordering follows the share of reads spent on the near-dead
container: 16\% for round-robin, 23\% for Thompson sampling, 25--27\% for tail-aware, UCB and
coverage-greedy, and 40--41\% for pivotal and learn-then-commit.  Information-greedy, which reads the
container with the largest posterior variance and spends 66\% of its reads on the near-dead container,
is 441\% costlier; its results are in Appendix~\ref{supp:info}.  With DT4DDS's
photolithographic synthesis model, no policy and none of the four static plans recovered the file in any
of 660 retrievals (eight policies at 10 syntheses $\times$ 3 read orders, four static plans, and
round-robin re-run at the full 30 $\times$ 10 design): every retrieval read the whole 72k-read stream and
reported failure, with no wrong output.  This result holds under that synthesis model with this code,
this decoder, a 24-read consensus cap and a 72k-read stream; it is a limit of our archive design and
decoder configuration, not evidence that photolithographic DNA archives are unrecoverable.

\emph{Model-free baselines.}\label{sec:res-extra}  Two simple rules that use no container model were
specified before they were run and were not tuned (Table~\ref{tab:r9base}(b)).  \emph{Successive
elimination} reads the active containers in turn and eliminates container $j$ once the upper end of the
interval (mean $\pm$ two standard errors, after at least three batches) on its new-fragment yield per
batch falls below the lower end of the best container's; if every active container returned no new
fragment in its last batch, all eliminated containers are reinstated.  The \emph{yield threshold} rule
reads the least-read container among those whose yield over their last four batches is at least 25\% of
the best container's, so a weak container is deprioritised but read again as soon as the others' yield
falls.  Both were run with round-robin on regenerated streams of the baseline, homogeneous and
one-failed scenarios.  On the regenerated baseline streams round-robin used 3.52k reads, against 3.49k on the committed streams.  On the baseline, successive elimination never eliminated a container and reproduced round-robin exactly, and the yield threshold gave $-$1.0\% ($-$1.6 to $-$0.6\%): the near-dead container's early yield is too close to the others' for either rule to set it aside.  In the homogeneous control both rules reproduced round-robin exactly.  With one failed container, successive elimination gave $-$0.5\% ($-$0.9 to $-$0.2\%) and the yield threshold gave $-$3.8\% ($-$5.5 to $-$2.4\%) on the retrievals that succeeded.  On these regenerated streams 50\% of one-failed retrievals succeeded, against 34\% on the committed streams, a reminder that recoverability after a failure depends on which strands each synthesis loses.  None of the 3{,}600 retrievals on regenerated streams returned a wrong file.

\subsection{Why adaptive allocation loses, and the dropout repair}\label{sec:res-misspec}\label{sec:res-dropout}\label{sec:res-calib}
The primary controller models a container by one efficiency multiplier $\theta_j$ that scales every
fragment's coverage.  The near-dead container (two half-lives, 40 PCR cycles, $\sigma{=}0.5$, three
copies per strand) instead loses a large fraction of its molecules outright: many of its fragments are
never read however long it is sequenced, while the surviving ones are read often.  Early batches from it
therefore recover fragments at a rate the model explains with a moderate $\theta_j$ and extrapolates to
the missing fragments, so coverage-based scores keep predicting gains there.  The posterior reflects
this misspecification: at $\rho_U{=}0.1$ its 90\% credible interval covered the target efficiency of the
five healthy containers in every held-out replicate but that of the near-dead container in only 0.20 of
cases, and no value of $\rho_U$ in the extended grid covers the near-dead container
(Appendix~\ref{supp:calib}).  The controller is therefore not calibrated where it matters,
and we describe it as online-adaptive rather than self-calibrating.

The ablation (Table~\ref{tab:r4abl}) separates three candidate causes: the container model, the
likelihood ($\rho_U$) and the prior width ($s$).  No single factor is \emph{the} cause.  For
coverage-greedy, paired against round-robin on the same 300 syntheses and read orders, the primary
corner costs $+$15.1\% (95\% CI $+$10.3 to $+$20.8\%); the binomial likelihood alone gives $-$4.4\% (95\% CI $-$6.0 to $-$2.8\%); widening the prior
alone gives $-$0.6\% (95\% CI $-$2.4 to $+$1.5\%); and adding dropout alone gives $+$1.6\% (95\% CI $-$0.1 to $+$3.3\%).  Thompson sampling behaves the
same way ($+$7.4\%, then $-$3.5\%, $+$0.9\%, $+$0.7\%).  The common mechanism is how fast the controller
can conclude that a container is bad: the over-dispersed likelihood widens the posterior, a tight prior
holds it near $\theta{=}1$, and a single-multiplier model keeps extrapolating the surviving fragments'
rate to the absent ones.  Relaxing any one of the three lets the evidence accumulate fast enough.  The
reduction is not uniform across rules: tail-aware keeps $+$10\% with the dropout model alone, and
pivotal and learn-then-commit keep $+$7\% to $+$13\% under single relaxations.

The dropout repair is the physically meaningful one.  A joint posterior over $(\theta_j,d_j)$ with
$d_j\in\{0,0.05,\allowbreak0.15,\allowbreak0.3,\allowbreak0.5,\allowbreak0.7,\allowbreak0.9\}$ attributes a fragment still missing after many reads increasingly
to permanent absence rather than to low coverage.  On the baseline, keeping the primary likelihood and
prior and changing only the model, UCB goes from $+$15.3\% to $-$3.5\% against round-robin (95\% CI
$-$4.8 to $-$2.2\%), coverage-greedy from $+$15.1\% to $+$1.6\%, Thompson sampling from $+$7.4\% to $+$0.7\%,
pivotal from $+$49.3\% to $+$7.4\% and learn-then-commit from $+$44.6\% to $+$12.7\%.  Coverage-greedy's
share of reads on the near-dead container falls from 0.27 to 0.20 and UCB's from 0.27 to 0.17, against
round-robin's 0.16.  In the homogeneous control, where there is no dropout to detect, the two models agree
within 0.5\% on every rule, so the second parameter costs nothing when it is not needed.  The clearest
benefits are in the scenarios the one-parameter model misreads most.  After an outright failure the
conditional saving grows from $-$17.0\% to $-$35.0\% for coverage-greedy, from $-$9.2\% to $-$36.2\% for
Thompson sampling and from $+$1.6\% to $-$38.5\% for tail-aware, with the success probability unchanged at
0.34; adding the binomial likelihood as well takes all four rules to between $-$42\% and $-$45\%.  Under
partial physical loss and common-cause degradation the two-parameter model turns penalties of $+$77\% and
$+$47\% into savings of $-$6\% and $-$34\% (Section~\ref{sec:res-robust}).

The dropout model is not uniquely responsible for the repair (the binomial likelihood alone does as much
on the baseline), it does not move the tail (Section~\ref{sec:res-tail}), and on the baseline its best
result, UCB's $-$3.5\%, is about what fixed proportions give rather than an addition to it.  Allocation
reaches 6--45\% only in dropout-dominated scenarios, where it repairs a problem the primary model created.

\subsection{Tail metrics and the static family}\label{sec:res-tail}
\input{tables/r4_tails}
No adaptive rule lowered the 0.99-quantile \emph{reliably}; fixed proportions did.  A q99 estimated from
300 retrievals is driven by its three most expensive observations, which come from a few synthesis
replicates, so we report cluster-bootstrap intervals over syntheses and re-ran the baseline with 100
independent syntheses $\times$ 3 read orders (Table~\ref{tab:r4tails}).  We distinguish a numerical
reduction in the estimated q99 from a reduction whose interval excludes zero, and quote the paired
differences of the table's $\Delta$q99 column.

In the primary design (30 syntheses $\times$ 10 orders), the pre-specified fixed proportions
($\alpha{=}1$) have q99 4.80k against round-robin's 5.15k, a paired difference of $-350$ reads (95\% CI $-450$ to $-3$),
or $-$6.8\%; the interval that excludes zero only narrowly.  In the 100-synthesis design they
reach 4.70k against 5.15k, $-450$ reads (95\% CI $-550$ to $-202$), or $-$8.7\%.  Three rules with the dropout model
are numerically below round-robin in at least one design: UCB ($-450$ reads, 95\% CI $-1{,}001$ to $+244$) and coverage-greedy
($-52$ reads, 95\% CI $-602$ to $+350$) in the primary design, and coverage-greedy ($-100$ reads, 95\% CI $-599$ to $+197$) and Thompson sampling
($-49$ reads, 95\% CI $-500$ to $+305$) in the 100-synthesis design.  Every one of these intervals contains zero.  Every
primary-model rule is worse than round-robin at q99 in both designs, with intervals that exclude zero,
including the tail-aware rule that was designed to target the bottleneck stripes, so that heuristic did
not achieve its aim.  Adaptive rules therefore did not reliably lower the tail, although we cannot rule
out that the best of them lowers it by up to about 9\%.  We report q95 beside q99 because it is the
better-determined tail statistic at this sample size.

The static family's members are $w_j\propto e_j^{\alpha}$, and the pre-registered grid was
$\alpha\in\{0,0.5,1,2\}$, so ``best static plan'' means the best of those four, not the optimal fixed
allocation.  A user who must reserve capacity gains nothing from stopping unless unused capacity is
refunded; for that user the relevant comparison is between reservations, and fixed proportions tuned to
the containers are the better choice.  In the homogeneous control the static plans must reserve 2.85k
against a round-robin consumed mean of 2.26k.  With a failed container no static plan can be sized at
all: it either reads the whole stream or fails, so its reservation is the full 72k reads.

\subsection{Secondary evidence: real read streams}\label{sec:res-replay}
\input{tables/replay}
Table~\ref{tab:replay} gives the replay with the generic prior.  On the Twist pool, round-robin with
verified stopping used 58.6k reads on average against the 65.7k budget of the hindsight-tuned static plan
($-$11\%), and the model-based rules used 56.6--57.2k, 2--3\% fewer than round-robin; coverage-greedy was
9.0k (95\% CI 8.5--19.3k) below the hindsight static budget.  The calibrated belief changed
coverage-greedy by $-$0.2\%.  As on the digital twin, no \emph{adaptive} rule lowered q99 materially:
65.6--67.2k for the model-based rules against the 65.7k static budget.  Stopping supplies most of the
gain (about 7.1k of 9.0k reads).  On the Genscript pool, whose PCR-30-cycle container is near-dead,
verified stopping with round-robin used 637k reads against 852k for fixed depth ($-$25\%), but 1\% more
than the 631k of the hindsight-tuned static plan, which knows in advance to give that container 2.7\% of
its reads (paired difference $+$6.0k, 95\% CI $-$57.4 to $+$15.1k).  Every model-based rule with the
generic prior was 0--3\% costlier than round-robin (UCB 637k, coverage-greedy and pivotal 639k, tail-aware
644k, learn-then-commit 651k, Thompson sampling 654k), and, as on the digital twin, each gave the
near-dead container more reads than round-robin did (19--20\% against 17\%).  The exception is the
calibrated belief, fitted on calibration reads of the same containers that are disjoint from the
evaluation reads: coverage-greedy with it gave the near-dead container 3.4\% of its reads and used 533k
($-$16\% vs.\ round-robin; $-$97.9k, 95\% CI $-$161.6 to $-$90.0k, vs.\ the hindsight static budget).
Allocation therefore pays when the weak container is identified \emph{before} retrieval, not when the
controller must learn it online.  q99 was 820--853k for the generic-prior rules and 676k with the
calibrated belief, against the 631k static budget.  Information-greedy did not complete on this pool and
is treated as censored (Appendix~\ref{app:devfull}).

\input{tables/r4_calcost}
\emph{What the calibrated prior costs.}  That 16\% is the largest allocation effect in this paper, and
it is not free: the calibrated belief is fitted on the first 300{,}000 reads of each of the six runs,
1.8M reads in total, which is 2.8 times what one retrieval costs.  Table~\ref{tab:r4calcost} charges
them.  Beliefs built from 1k, 3k or 10k reads per container (counts only, generic consensus curves)
save nothing measurable (640k, 638k and 636k reads against round-robin's 637k), and 30k per container is
worse (658k), because a half-informative estimate of a container's abundance distribution misleads the
allocator more than no estimate at all.  Only at 100k reads per container (577k, 600k calibration reads)
and at the full 300k (533k, 1.8M calibration reads) does the saving appear.  Dividing the calibration
reads by the per-retrieval saving gives the break-even point: about 10 retrievals for the 100k belief and
17 for the full one.  Calibration therefore pays only for an archive that will be read repeatedly, and for
a single retrieval it costs several times what it saves.

\subsection{Robustness}\label{sec:res-robust}\label{sec:res-fail}\label{sec:res-uncond}
\emph{Failed containers.}  With one failed container (the fresh one, which returns only non-informative
reads), 34\% of retrievals succeeded under every policy, with identical success on every synthesis and
read order: recoverability is set by the archive's redundancy, not by the policy.  On the 102 retrievals
where all policies succeeded, adaptive allocation helped, because the failed container is easy to
identify and stops being read: UCB used 24\% (17--28\%) fewer reads than round-robin, coverage-greedy
17\% (9--23\%), pivotal 15\% (6--21\%) and Thompson sampling 9\% (1--15\%); tail-aware was even ($+$2\%, $-$7 to $+$14\%) and
learn-then-commit 8\% (3--14\%) costlier (Table~\ref{tab:dt}).  Counting all 300 attempts, including the
66\% that read the whole stream and failed, UCB's saving is 3\% (Table~\ref{tab:uncond}).
With two failed containers no policy recovered the file.  Both results are limits of the archive design.
With round-robin placement of $(32,48)$ stripes over six containers, each failed container removes 8 of a
stripe's 48 fragments, half of its redundancy of 16.  After one failure a stripe tolerates only 8 further
lost strands: 18 of the 30 syntheses were unrecoverable under every read order, 7 always recoverable, and
5 depended on the read order; after two failures exactly $K{=}32$ fragments remain, so a single strand
lost by synthesis, aging or PCR in any stripe makes the file unrecoverable.  Surviving more failures
requires more redundancy or a different placement; no read-allocation policy can compensate.

\emph{Severity, prior and decoder.}  Varying only the near-dead container's copy number
(Fig.~\ref{fig:sweep}), adaptive rules are 2--3\% cheaper than round-robin at the mild end
(200 and 30 copies) and costlier at 3 and 1 copies, so adaptive allocation helps only while the model's
single multiplier is a fair description of the worst container.  Varying the controller's prior
assumptions one at a time leaves coverage-greedy between $+$10.9\% and $+$20.6\% against round-robin
under the one-parameter model, except for a tighter prior width ($+$56.4\%), and between $-$0.3\% and
$+$3.4\% under the two-parameter model (Appendix~\ref{supp:robust}).  The consensus cap
(reads per cluster used to form a consensus, 24) is not binding on the baseline, where caps of 12 and 48
give the same means to within 10 reads, but it binds where a stream is read to exhaustion: with one
failed container, round-robin recovers the file in 31.3\% of retrievals at a cap of 12, 34.0\% at 24 and
35.7\% at 48, with no wrong output at any setting.  Part of the redundancy limit above is therefore a
limit of our decoder configuration, and every zero- or partial-success result in this paper holds only
for this code, placement, consensus cap and stream length.

\emph{Failure modes and geometries.}  Four further failure modes (Table~\ref{tab:r4fail})
confirm the dropout explanation: under partial physical loss of half of one container's designs, the
one-parameter model costs $+$76.6\% (coverage-greedy), $+$92.0\% (tail-aware) and $+$59.7\% (UCB), while
coverage-greedy with dropout reaches $-$6.0\%, and $-$33.6\% under common-cause degradation.  Success
probability again does not depend on the policy, and no configuration produced a wrong output.  Six
further geometries (Table~\ref{tab:r4geom}: 3 and 12 containers, $(32,40)$, $(32,64)$ and
$(16,24)$ codes, and a file four times larger) preserve every qualitative conclusion.  Coverage-greedy
under the primary model is never cheaper than round-robin ($+$1.1\% to $+$93.5\%); with the dropout model
it ranges from $-$27.9\% to $+$2.4\%; the hindsight static plan with stopping is cheaper than round-robin
everywhere ($-$4.0\% to $-$10.7\%); and there was no wrong output in any geometry.

\subsection{Reads are not the only cost}\label{sec:res-cost}
\input{tables/r4_cost}
Our primary metric is reads, which is the right unit for a pay-per-read service and the wrong one for a
laboratory that pays per run.  Table~\ref{tab:r4cost} reprices the same 300 baseline retrievals as
Tables~\ref{tab:actual} and~\ref{tab:dt} under three charging models: reads only; a multiplexed run that
charges 50 read-equivalents per batch of 50 reads and 500 per container library preparation; and a
separate sequencing run per batch at 2{,}000 read-equivalents.  The ranking of the policies is unchanged,
because every policy's batch count is proportional to its read count and every policy prepares all six
libraries.  What changes is the margin: under the per-batch-run model the whole spread between
round-robin and the best adaptive rule is a few per cent of a bill dominated by run setup.  Container
switches per retrieval were recorded only by part of the runs.  On these 300 retrievals, round-robin
switches 63 times and the rules with the dropout model 38--59 times; on the 210 retrievals of
the 70-synthesis extension, the primary-model rules switch 38--61 times (tail-aware fewest), round-robin
63 times and the fixed-proportion plans 65 ($\alpha{=}1$) and 69 (uniform) times, because they cycle to
hold their proportions.  No policy here is switch-frugal by design; a setting where switching is
expensive would need a larger batch or an explicit switch penalty in the score.

Three quantities are distinct: reads \emph{consumed} (what a policy spends), capacity \emph{reserved}
(what a non-stopping plan must commit in advance), and laboratory \emph{expenditure} (what an invoice
says).  Our 27--32\% figure compares the first with the second; it is \emph{not} a demonstrated
reduction in the cost of sequencing, which would require per-read pricing with no fixed component.  We
did not model instrument time, reagent costs, library preparation, minimum run sizes or flow-cell
granularity.

%=====================================================================
\section{Discussion}\label{sec:disc}
%=====================================================================
\subsection{Why stopping gains more than allocation}\label{sec:disc-theory}
Write $T$ for the number of reads a retrieval needs under a fixed read schedule.  A plan that does not
stop must commit a budget $B$ with $\Prob(T\le B)\ge0.99$, so it pays the q99 of $T$ on every retrieval;
verified stopping pays $T$ itself, on average $\E[T]$.  The stopping saving is therefore
$1-\E[T]/\mathrm{q99}(T)$, which depends only on the spread of $T$ relative to its mean.  On the baseline,
$T$ has a coefficient of variation of 0.16 and its q99 lies 48\% above its mean, which is
the 32\% saving against the uniform reservation.  That spread has two sources.  Within one synthesis the
read order is random, which accounts for 69\% of the variance of round-robin's read count;
between syntheses, abundance and dropout differ, which accounts for the remaining 31\%.  A
reservation must cover both, whereas stopping pays only for the realised retrieval.

Allocation acts on a different quantity, the number of new verified fragments per read.  Consider a fluid
approximation in which container $j$ yields new verified fragments at a constant rate $e_j$ per read and,
because the $K$-of-$N$ code makes containers substitutable, any fragment counts towards completion.
Round-robin then collects fragments at the average rate $\bar e$, and no allocation can exceed the best
rate $e_{\max}$, so under this approximation the saving of any allocation over round-robin is at most
$1-\bar e/e_{\max}$.  With the per-container rates measured on the baseline (0.065--0.124 verified
fragments per read), this bound is 11\%.  Concentrating reads on the best containers also
exhausts their recoverable fragments sooner, so the attainable gain is smaller; the best fixed plans
reach 4.2--4.6\%.  The bound depends on the spread of the container rates, not on how well a controller
learns them, which is why even a plan that knows the rates in hindsight gains little, and why a
controller that misjudges one container can easily lose more than the bound allows it to win.  The
argument is approximate: it ignores the stripe structure and saturation, and it is not an optimality or
regret result.

\subsection{What changes at larger scale}\label{sec:disc-scale}
$T$ is the maximum over the $S$ stripes of the read counts at which each stripe completes.  If these were
independent, the mean of the maximum would grow roughly logarithmically in $S$ while its spread stayed of
the same order, so the read-order part of the stopping saving would shrink relative to the mean as the
archive grows.  The synthesis-level part does not average out, because a weak or failed container affects
every stripe.  The geometries we ran are consistent with this.  With a file four times larger (100
stripes), the stopping saving against the uniform reservation falls from 32\% to 24\%
(coefficient of variation 0.14); with the higher-redundancy $(32,64)$ code it is 18\%,
and with 12 containers it rises to 41\%.  These geometries have only 30--60 retrievals, so
their q99 is close to the sample maximum and imprecise.  At 100 times more strands we would expect a
smaller relative stopping saving, increasingly set by synthesis-level variation and container failures,
which a reservation must cover and stopping does not; we did not run that scale.

\subsection{What helps, and whom}
The robust gain in our experiments is verified stopping: reading until every stripe, its CRC-32 and the
file hash verify, and then stopping.  It needs no model of the containers, never returned a wrong file,
and for a user who pays per read it consumes 27--32\% fewer reads than a non-stopping plan must reserve.
It does \emph{not} reduce the capacity a user must reserve, because the q99 is unchanged.  Choosing
\emph{which} container to read next is a second-order effect on the baseline: worth a few per cent at
best, and strongly negative when the container model cannot express what is wrong with a container.  The
two best cases reported here come from \emph{different experiments} and should not be added: the 4.2\%
saving is the pre-specified fixed allocation, which runs no controller and whose proportions can be
estimated offline; the 3.5\% saving is UCB, an online rule, after the container model is extended with a
dropout parameter.  They are two routes to a similar small gain, not a cumulative 7.7\%.  Read savings are
also not laboratory-cost savings, and redundancy must be sized for the failures the archive has to
survive: with $(32,48)$ over six containers, two lost containers leave no margin at all.

\textbf{Calibration versus cost.}  The likelihood that gave the best-calibrated intervals
($\rho_U{=}0.1$) also gave the costlier allocations, and the binomial likelihood, whose intervals covered
only 54\%, gave cheaper ones.  Calibrating the uncertainty of a misspecified model does not improve its
decisions.  No likelihood in the family covers the near-dead container, a failure of the container model
that only a dropout parameter addresses; but on cost the likelihood matters a great deal, since at
$\rho_U{=}0.03$ most of the allocation penalty disappears without changing the model
(Section~\ref{sec:res-alloc}).  A configuration can therefore be poorly calibrated and still allocate
acceptably.  The practical priority is to repair the container model before tuning the posterior over it.

\textbf{Digital twin versus replay.}  The two settings agree that stopping is the main saving and that
no adaptive rule lowers q99 reliably.  On the Twist pool, whose containers are runs of one synthesis
batch with no near-dead member, allocation saved 2--3\%; on the Genscript pool, which has a near-dead
container, it cost 0--3\% and over-read that container, as on the digital twin.  Only prior knowledge of
the weak container, obtained from calibration reads, turned allocation into a large saving.

\begin{figure}[t]
\centering
\fbox{\parbox{0.94\columnwidth}{\small
\textbf{Practical recommendation (multi-container retrieval).}
\begin{enumerate}
\setlength{\itemsep}{1pt}
\item \textbf{Stop on verification.}  Read the containers in turn and stop when every stripe matches its
CRC-32 and the file its hash.  This needs no container model and consumed 27--32\% fewer reads than a
non-stopping plan must reserve; testing every 50 reads loses less than 1\% against testing every read.
Never return an unverified file.
\item \textbf{Estimate fixed read proportions offline.}  Proportions measured once, in hindsight or on
disjoint calibration replicates, gave 4.2--4.4\% fewer reads than round-robin and are the only allocation that lowered
the 0.99-quantile reliably ($-$6.8\% in the primary design, $-$8.7\% in the
100-synthesis design), which is what matters to a user who reserves capacity.
\item \textbf{Calibrate only for archives that will be read repeatedly.}  A prior fitted on held-out reads
of the same containers saved 16\% per retrieval on the Genscript replay but consumed 1.8M calibration
reads, 2.8 times one retrieval, and broke even only after about 17 retrievals; smaller calibration budgets
did not help consistently.
\item \textbf{If you must adapt online, use a dropout-aware container model, and deprioritise rather than
exclude weak containers.}  Under the pre-registered likelihood, rules with a single-efficiency model read
7--49\% more than round-robin; with efficiency and dropout the best online rule gained 3.5\%.  A container
that yields nothing in one probe may be low-efficiency, partly lost or in need of deeper consensus, so
keep sampling it.
\end{enumerate}}}
\caption{Practical recommendation derived from Section~\ref{sec:results}.}
\label{fig:recommend}
\end{figure}
Fig.~\ref{fig:recommend} summarises the practical recommendations that our results support for
archives like ours.

\textbf{Future experimental validation.}\label{sec:wetlab}  Physically separate containers would allow
three tests.  (i) Whether real containers show the near-dead, dropout-dominated behaviour that made
adaptive allocation lose here: one design would be synthesised in several independent batches, aged and
amplified differently, and each container sequenced to saturation to measure its fraction of absent
strands.  (ii) Whether verified stopping keeps its saving and its safety on real sequencing runs with our
own inner code: reads from a nanopore or a split Illumina run would be streamed through the same decoder,
counting wrong outputs and inner-code false accepts against the known file.  (iii) What real container
failures look like and how often they occur, which would calibrate the failure prior and the redundancy
needed per container.  The decisive experiment, not performed here, is the full chain on physical
material: encode a known file, synthesise it into independently handled containers, sequence each one,
feed the FASTQ reads into this decoder, and compare the returned bytes with the original.

%=====================================================================
\section{Limitations and Threats to Validity}\label{sec:limits}
%=====================================================================
Table~\ref{tab:limits} lists the main threats.  \emph{External validity}: every result is
computational.  The digital twin uses DT4DDS's default parameters and one primary archive geometry
(1{,}200 strands, $(32,48)$, six containers), and its failure modes are simulated.  The replay uses real
reads, but decodability is an oracle event and the ``containers'' are runs of one synthesis batch.
\emph{Internal validity}: $\rho_U$ was selected with prior width $s{=}1$ and used with $s{=}0.5$; the
selected value lies at the edge of the pre-specified grid; one code fix (static allocator) was made after
the protocol was fixed; and the per-read-stop and model-free baselines run on regenerated read streams,
so they are paired only with the round-robin re-run on those streams.  The pre-specified static plans use
efficiencies measured in hindsight on the evaluated syntheses, which favours them; the
calibration-estimated variant, which sees no evaluation data, performs the same within noise.
\emph{Construct validity}: read counts are a proxy for cost.  \emph{Statistical}: wrong-output and
false-accept bounds come from zero events and scale with the number of trials; the comparisons are
paired but within one primary archive design; intervals are not adjusted for multiplicity
(Section~\ref{sec:eval-stats}).  The analysis of Section~\ref{sec:disc-theory} is a fluid approximation;
we prove no optimality or regret result.
\input{tables/limitations}

%=====================================================================
\section{Reproducibility and Availability}\label{sec:repro}
%=====================================================================
\textbf{Artifact.}  Code, configuration, every result file and the scripts that build every table and
figure are in the public repository \anon[(link withheld for anonymous review)]{\url{https://github.com/SaiMouli3/Dev}} (directory
\texttt{d3ra-dna}), at the commit recorded in \texttt{research\_paper/final/COMMIT.txt}.  The environment
is Python 3.13.18 with NumPy 2.5.3, SciPy 1.18.1, pandas 3.0.5 and DT4DDS~1.1.0 (pinned in
\texttt{requirements.txt}); the public data are ENA PRJEB65931~\cite{gimpel2023prjeb65931} (Illumina) and
the DNAformer nanopore reads~\cite{technion2025dnaformerreads,barlev2024binned}, both fetched by the scripts
in \texttt{experiments/}.  The digital-twin experiments are driven by
\texttt{experiments/dt\_experiments.py}, the replay by \texttt{experiments/paper\_experiments.py};
statistics by \texttt{dt\_analysis.py}, \texttt{r4\_analysis.py} and \texttt{r9\_analysis.py}, which share
one cluster-bootstrap module; tables and figures by \texttt{make\_tables.py}, \texttt{make\_r4\_tables.py},
\texttt{make\_r9\_tables.py} and \texttt{make\_dt\_figures.py}, which read only committed result files.
\texttt{research\_paper/audit/} recomputes every quantitative claim from the result files and checks it
against the PDF.  Table~\ref{tab:evidence} maps each claim to its command and result file.

\textbf{Determinism.}  Our own layer is deterministic: every random draw in placement, read ordering, the
controller and the decoder is seeded by CRC-32 of a string naming the scenario, replicate, container, read
order and policy, and 24 sampled baseline retrievals re-executed single-process in reverse order reproduce
every recorded field exactly.  The \emph{simulator} is not: DT4DDS~1.1.0 does not produce identical read
streams from a fixed seed.  The committed read streams are therefore part of the artifact, and a reader
who regenerates them does not get our retrievals back.  Regenerating all 30 baseline replicates from a
clean clone and re-running all eight policies and four static plans (3{,}600 retrievals) left every
outcome identical (the same successes, no wrong outputs) while 2{,}350 of 2{,}400 retrievals consumed a
different number of reads (mean absolute difference 1{,}629).  Every policy's mean fell inside the other
run's cluster-bootstrap interval in both directions, the ordering of all twelve policies was identical,
and the headline comparisons moved by one to three percentage points (round-robin 3.49k against 3.51k;
the hindsight static plan $-$4.2\% against $-$5.4\%; coverage-greedy $+$15.1\% against $+$12.2\%).  The
comparison is in \texttt{evidence/}\allowbreak\texttt{verification\_reports/}\allowbreak\texttt{clean\_repro.txt}.

%=====================================================================
\section{Conclusion}
%=====================================================================
On a digital twin whose containers are independently synthesised \emph{in simulation}, verified stopping
returned the stored file without a single wrong output in \SafeN{} retrievals and consumed 27--32\% fewer
reads than the capacity a non-stopping plan must reserve.  Smaller reservation margins fail often, and
testing after every read instead of every batch adds less than one per cent.  This saving requires no
model of the containers and is the main result of the paper.

Deciding \emph{which} container to read next is worth far less.  Fixed proportions, which run no
controller and can be estimated offline, gave 4.2\% on the baseline and were the only allocation that
lowered the 0.99-quantile reliably.  Under the pre-registered single-efficiency model and likelihood,
model-based rules cost 7--49\% more than round-robin when a container was near-dead; under the likelihood
chosen by a corrected selection most were indistinguishable from round-robin, and relaxing any one of
model, likelihood or prior width removed most of the penalty.  On the baseline no adaptive rule, and
neither of two model-free rules, beat round-robin by more than a few per cent.  Larger allocation gains
appeared only in dropout-dominated scenarios, up to $-$45\% after an outright failure, and with a prior
calibrated on held-out reads ($-$16\%), which pays for itself only after about 17 retrievals.  A fluid
approximation explains the asymmetry: stopping recovers the spread of the read requirement, whereas
allocation is bounded by the spread of the containers' recovery rates.

The digital-twin experiments run the full decode chain and compare bytes; the real-read replay uses
public Illumina streams but treats decodability as an oracle event and emulates containers by subsetting
runs of one synthesis batch.  No~DNA~was synthesised or sequenced for this paper, and each headline figure
is a read count, not a laboratory cost.  The open questions require physically independent containers with
measured failure and dropout statistics, and an end-to-end retrieval of a known file from synthesised DNA
through this decoder.

\appendix

%=====================================================================
\section{Positioning against prior work}\label{supp:lit}
%=====================================================================
Table~\ref{tab:literature_comparison} compares this work with the closest prior work discussed in
Section~\ref{sec:related}.
\input{tables/literature_comparison}

%=====================================================================
\section{Operational specification of the retrieval loop}\label{supp:alg}
%=====================================================================
Table~\ref{tab:algorithms} specifies the loop and every allocation rule: initialisation, batch size,
tie-breaking, when the posterior is updated, the stopping test and how an exhausted stream is handled.
The two model-free baselines of Table~\ref{tab:r9base}(b) are specified in Section~\ref{sec:res-extra}.
\input{tables/algorithms}

%=====================================================================
\section{Information-greedy allocation}\label{supp:info}
%=====================================================================
Information-greedy reads the container with the largest posterior variance of $\ln\theta_j$ until every
variance falls below $10^{-3}$, then switches to coverage-greedy.  Under the primary container model the
near-dead container's variance never falls below the threshold, so the rule spends 66\% of its reads
there, against 16\% for round-robin, and costs 441\% more than round-robin on the baseline.  A richer
model gives it more uncertainty to chase: with the dropout model the penalty grows to 612\%.  In the
homogeneous control it is indistinguishable from round-robin, and after an outright failure it is 81\%
costlier on the retrievals that succeed.  Table~\ref{tab:suppinfo} collects every result for this rule.
It is a pure information-gathering rule with no cost term, and we report it for completeness rather than
as a competitive baseline.
\input{tables/supp_infogreedy}

%=====================================================================
\section{One failed container: conditional and unconditional cost}\label{supp:uncond}
%=====================================================================
\input{tables/r4_uncond}
With one failed container 34\% of retrievals succeeded, and the success probability was identical for
every policy and every static plan (Table~\ref{tab:uncond}): whether the file can be recovered is decided
by which fragments survive in the five remaining containers, not by the reading order.  Among the 102
retrievals that succeeded, UCB used 15.49k reads against round-robin's 20.31k, 24\% fewer.  Counting all
300 attempts, including the 66\% that read the whole 72k-read stream and failed, UCB used 52.79k against
round-robin's 54.42k, a 3\% reduction.  The two numbers answer different questions: a user who retries
until success pays the conditional number, and a user who pays per attempt pays the unconditional one.
After an outright failure, posterior rules therefore cut the cost of the retrievals that succeed by up to
24\% without improving the probability of success, and change the cost per attempt by 3\%.

%=====================================================================
\section{Calibration of the posterior and re-selection of $\rho_U$}\label{supp:calib}
%=====================================================================
\input{tables/r4_calib}
On the ten selection replicates, the 90\% credible-interval coverage of $\theta^\ast$ rose with $\rho_U$
from 0.45 (binomial) through 0.63, 0.73 and 0.80 to 0.87 at $\rho_U{=}0.1$, the edge of the pre-specified
grid, which was selected.  On the 30 held-out replicates it gave 0.87 overall (0.84 at the earlier and
0.89 at the later read count), 1.00 for five containers and 0.20 for the near-dead one
(Fig.~\ref{fig:calib}).  That selection has two defects: it used prior width $s{=}1$ while every policy
run uses $s{=}0.5$, and the selected value lies at the grid boundary.

A secondary analysis repeats the selection with $s{=}0.5$ and the extended grid
$\{0,0.003,\allowbreak0.01,\allowbreak0.03,\allowbreak0.1,\allowbreak0.2,\allowbreak0.3,\allowbreak0.5\}$ (Table~\ref{tab:r4calib}).  It re-runs the selection procedure
only.  Under the prior actually used, the best value is $\rho_U{=}0.03$ (selection coverage 0.81, held
out 0.83) rather than 0.1 (0.81 and 0.82), a difference within noise.  Extending the grid does not
rescue calibration: coverage flattens near 0.80 for $\rho_U\in[0.03,0.3]$ and falls to 0.60 at 0.5, so
no value in this likelihood family reaches the 0.90 target.  At every $\rho_U$ the five healthy
containers are covered 0.92--1.00 of the time and the near-dead container 0.00--0.03 of the time.
Widening a likelihood cannot fix a container whose behaviour the model cannot express.  The complete
policy comparison at the re-selected value is reported in Section~\ref{sec:res-alloc} 
(Table~\ref{tab:r6lik}).
\begin{figure*}[t]
\centering
\includegraphics[width=0.98\textwidth]{fig_dt_calibration.pdf}
\caption{Is the posterior over container efficiency calibrated?  Each point or bar is the fraction of
cases in which the 90\% credible interval for $\theta_j$ contained the target $\theta^\ast_j$; a
calibrated posterior would sit on the dotted line at 0.90.  \textbf{(a)} Selection replicates
1000--1009, one curve per container, against the beta-binomial correlation $\rho_U$; the thick black
curve is the mean over containers, which is what the selection rule maximised.  \textbf{(b)} Held-out
replicates 0--29 at the selected $\rho_U{=}0.1$, per container.  Five containers are covered perfectly
and the near-dead one is covered in 0.20 of cases at this $\rho_U$ (0.00--0.03 across the extended grid
of Table~\ref{tab:r4calib}).}
\label{fig:calib}
\end{figure*}

%=====================================================================
\section{Completion forecast}\label{supp:forecast}
%=====================================================================
The controller also forecasts the probability of finishing within a remaining budget.  The
\emph{analytic} forecast multiplies per-stripe normal approximations of
$\Prob(\sum_{i\in s}X_i\ge\nu_s)$ with the remaining budget spread evenly over containers.  The
\emph{posterior-simulation} forecast draws 200 values of $\theta$ from the posterior, spreads the
remaining budget in proportion to the reads spent so far, simulates per-fragment recoveries and reports
the fraction of draws in which every stripe completes.  It was evaluated with a budget equal to the
median coverage-greedy cost, at 30/50/70/85\% of it, by expected calibration error (ECE, 10 bins), Brier
score against a constant (climatology) forecast~\cite{gneiting2007strictly} and AUC.

No policy depends on the forecast; we report it because it was pre-specified.  On 1{,}127 checkpoints
from 300 baseline retrievals with a budget of 3.75k reads, of which 52\% finished within the budget, both
forecasts were over-confident (mean forecast 0.86 against 0.52 observed at 30\% of the budget).  The
posterior-simulation forecast beat the analytic one on every metric (ECE 0.29 vs.\ 0.34; Brier 0.30, 95\%
CI 0.27--0.34, vs.\ 0.34, 0.29--0.38; AUC 0.73, 0.69--0.78, vs.\ 0.69, 0.65--0.73), yet both Brier scores
are worse than the 0.25 of a constant forecast at the observed base rate, so neither beats climatology.
They rank retrievals usefully only late in a retrieval (AUC 0.82 at 85\% of the budget).  Futility
stopping, which would rely on them, is disabled in every run, and the forecasts should not be used as
probabilities.
\begin{figure}[t]
\centering
\includegraphics[width=0.6\columnwidth]{fig_forecast.pdf}
\caption{Reliability of the completion forecast on the digital twin: forecast probability of finishing
within the remaining budget against the observed frequency of finishing, pooled over checkpoints at
30/50/70/85\% of the budget (10 equal-width bins, bins with fewer than five checkpoints dropped).  A
calibrated forecast would lie on the diagonal; both curves lie below it, so both are over-confident, and
both have a worse Brier score than a constant forecast at the base rate.}
\label{fig:forecast}
\end{figure}

%=====================================================================
\section{Severity sweep}\label{supp:sweep}
%=====================================================================
The severity sweep (Fig.~\ref{fig:sweep}) varies only the near-dead container's copies per strand.  At
200 and 30 copies, coverage-greedy and Thompson sampling are 2--3\% cheaper than round-robin (2.65k vs.\
2.72k at 200 copies) and tail-aware is even; at 10 copies the policies are within $\pm$1\% except
tail-aware ($+$5\%); at 3 and 1 copies all adaptive rules are costlier (at one copy: coverage-greedy
4.69k, Thompson 4.51k, tail-aware 4.47k vs.\ 4.16k).  Success was 1.00 throughout.  The point at 3 copies
reproduces the baseline design with other seeds (round-robin 3.46k vs.\ 3.49k).  Adaptive allocation
therefore helps only while the model's single multiplier is a fair description of the worst container.
\begin{figure}[t]
\centering
\includegraphics[width=0.6\columnwidth]{fig_dt_sweep.pdf}
\caption{Severity sweep: mean reads consumed as the degraded container's copy number falls from 200 to 1
(300 retrievals per point and policy, 95\% cluster-bootstrap intervals; the axis runs from mild to severe
degradation).  Every retrieval at every level returned the correct file, so the success probability is
1.00 throughout and is not plotted.  Adaptive rules are cheaper than round-robin only at the mild end.}
\label{fig:sweep}
\end{figure}

%=====================================================================
\section{Robustness details}\label{supp:robust}
%=====================================================================
\textbf{Prior assumptions.}  The controller's generic assumptions were varied one at a time on the
baseline (30 syntheses $\times$ 3 read orders each): the prior probability of outright failure
($10^{-2}$, $10^{-1}$ against the primary $10^{-3}$), the prior dropout (0, 5\% against 0.5\%), the
negative-binomial shape (3, 30 against 10), the $\theta$ grid (40 points over $[0.001,10]$ and 10 points
over $[0.01,5]$ against 25 over $[0.01,5]$), the prior width ($s{=}0.25$, $1.0$ against $0.5$) and the
consensus-success curve (scaled by 0.5 and 2.0).  Under the one-parameter model, coverage-greedy stays
between $+$10.9\% and $+$20.6\% against round-robin in every setting except the tighter prior
$s{=}0.25$, where it rises to $+$56.4\%.  Changing the failure prior by two orders of magnitude has no
effect ($+$17.2\% either way), because the failure atom is not what is mismodelled.  Under the
two-parameter model the same sweep stays between $-$0.3\% and $+$3.4\%.

\textbf{Decoder configuration.}  Table~\ref{tab:r4cap} varies the consensus cap.
\input{tables/r4_cap}

\textbf{Failure modes beyond total loss.}  Table~\ref{tab:r4fail} gives the four additional failure
modes summarised in Section~\ref{sec:res-robust}.
\input{tables/r4_fail}

\textbf{Other archive geometries.}  Table~\ref{tab:r4geom} gives the six additional geometries.
\input{tables/r4_geoms}

%=====================================================================
\section{Supporting channel result (real nanopore reads)}\label{supp:nanopore}
%=====================================================================
On two independent nanopore flowcells of DNAformer reads~\cite{barlev2025deep,technion2025dnaformerreads},
a strand-specific error propensity persists across flowcells ($r{=}0.57$; null with matched contexts
$0.002$), and a hierarchical beta-binomial error model fitted on one flowcell predicts consensus success
on held-out strands of the other with mean absolute error 0.083, against 0.181 for an i.i.d.\ model.  This
motivates over-dispersed models but does not enter the controller.

%=====================================================================
\section{Upper bound with evaluation leakage}\label{app:leak}
%=====================================================================
Our first analysis of the real read streams (Table~\ref{tab:leak}) used a ``calibrated'' belief whose
per-strand consensus curves were pooled over reads that included the evaluation split.  The controller
therefore had information about the reads it was evaluated on, and the static plans were tuned on them
too.  We keep the table only as an upper bound with evaluation leakage; no claim in the main text uses
it.  In the analysis of Table~\ref{tab:replay} the primary belief is the generic prior, and the optional
calibrated belief uses only the first 300k reads of each run while evaluation uses the next 900k.  On
Twist, the leaky analysis put coverage-greedy at 55.7k reads and the oracle-tuned static plan at 63.1k;
the analysis without leakage gives 56.7k and 65.7k.  The leak therefore flattered the adaptive policy by
about 1k reads and the static baseline by about 2.6k, and the gap without leakage is larger, not smaller.
\input{tables/leak_appendix}

%=====================================================================
\section{Protocol deviations and post-hoc analyses}\label{app:devfull}\label{app:dev}
%=====================================================================
\textbf{Deviations.}  The protocol and every change made after it was fixed are recorded in
\texttt{dt\_protocol.yaml}.  (i) The planned CustomArray scenario was replaced, before any policy run, by
DT4DDS's photolithographic synthesis model, because the CustomArray settings carry a 117-nt
position-bias table that cannot be applied to 136-nt strands.  (ii) We suspected from commit timestamps
that the first Twist replay had run with $\rho_U{=}0$ and re-ran it; a direct test showed that it had
used $\rho_U{=}0.1$, and the re-run reproduced it exactly.  (iii) One cached homogeneous replicate was
truncated by an interrupted job; it was regenerated from its seed and the scenario re-run in full.
(iv) The photolithographic scenario was first run with 10 syntheses $\times$ 3 read orders for compute
reasons; the first 50 retrievals of the full run were discarded unseen.  Round-robin was later re-run on
it at the full 30 $\times$ 10 design, with the same outcome (0 of 300 recovered).  (v) The static
allocator crashed when only zero-weight (failed) containers remained; it now reads them out, which can
only affect failed retrievals of the failure scenarios.  The $\rho_U$ rule was applied as specified even
though the selected value lies at the edge of the grid.  (vi) The sequential Genscript replay finished
information-greedy after about six hours but then crashed while computing its paired difference, because
some trials never completed (infinite cost); nothing was saved.  We fixed the bug, ran the other Genscript
conditions in parallel processes with the same per-condition seeds (they reproduced every value logged by
the sequential run), and re-started information-greedy; it had not finished when this version was built.
In the sequential run its 200 trials finished and more than 1\% of them never recovered the file within
the available reads, which is why the 0.99 budget was infinite.  We treat the condition as
\emph{censored}: it is excluded from every average and comparison.  (vii) The secondary analyses of
\texttt{dt\_protocol\_rev4.yaml} were specified before they were run: the dropout-aware container model,
the $2\times2\times2$ ablation, the $\rho_U$ re-selection at $s{=}0.5$ on an extended grid, prior
sensitivity, four further failure modes, six further archive geometries, consensus caps of 12 and 48,
finer and calibration-estimated static families, an extended baseline replicate set, the full-design
photolithographic run, the replay calibration-cost experiment and the cost model.  (viii) The following analyses
are exploratory: the complete comparison at $\rho_U{=}0.03$, the reservation
baselines, and the per-read stop, successive-elimination and yield-threshold baselines on regenerated
streams (Section~\ref{sec:res-extra}).  (ix) Every cluster-bootstrap interval was
recomputed with one shared, order-independent resampling scheme (Section~\ref{sec:eval-stats}); point
estimates are unchanged, and interval limits moved by at most 0.8 percentage points on the baseline,
1.3 in the failure scenarios and 3.6 for information-greedy.

\textbf{Post-hoc analyses} (not pre-specified).  With the binomial likelihood ($\rho_U{=}0$), on the same
300 baseline retrievals, coverage-greedy used 3.34k reads ($-$4.4\% vs.\ round-robin, against $+$15.1\%
at $\rho_U{=}0.1$), Thompson sampling 3.37k ($-$3.5\%, against $+$7.4\%), tail-aware 3.37k ($-$3.5\%,
against $+$12.9\%) and pivotal 3.64k ($+$4.4\%, against $+$49.3\%); q99 was 4.95--5.70k, still not below
round-robin's 5.15k.  With one failed container the binomial likelihood helped tail-aware ($-$25\% vs.\
$+$2\%) and Thompson sampling ($-$20\% vs.\ $-$9\%) and hurt coverage-greedy and pivotal slightly
($-$15\% vs.\ $-$17\%; $-$8\% vs.\ $-$15\%).  Success and safety were unchanged (no wrong output).  On the
Twist replay, the binomial likelihood made coverage-greedy and Thompson sampling costlier (57.6k each,
against 56.7k and 56.9k at $\rho_U{=}0.1$), the opposite of the digital-twin result: the two settings
disagree on which likelihood allocates better.

%=====================================================================
\section{Claim-to-evidence map}\label{supp:evidence}
%=====================================================================
\input{tables/evidence}

\section*{Use of Generative AI}
In accordance with ACM's Policy on Authorship, we disclose that the generative AI tool Claude (Anthropic)
was used to assist with writing and with grammar and language editing of this manuscript.  The authors
reviewed and verified all content and take full responsibility for it.

\begin{acks}
We thank the authors of the public datasets and of DT4DDS (ENA PRJEB65931; DNAformer;
Srinivasavaradhan \emph{et al.}) for releasing reads, designs and software.
\end{acks}

\bibliographystyle{ACM-Reference-Format}
\bibliography{references}

\end{document}
